\chapter{Dimensionality Reduction}\label{n:dimred}
\studytime{2 hours}

\begin{dfn}{Dimensionality reduction}
\textbf{Dimensionality reduction} is the process of reducing the number of features (dimensions) of a dataset while keeping as much of the important information as possible.
\end{dfn}
\textbf{Why:} speed up training, remove noise and redundant features, avoid the curse of dimensionality, compress data, visualise data in 2-D or 3-D (\textbf{data visualisation}). \textbf{Cost:} some information is lost; pipelines become more complex.

\section{The Curse of Dimensionality}
\begin{dfn}{Curse of dimensionality}
The \textbf{curse of dimensionality} refers to the problems that arise when data has very many features: the space becomes so large that data is \textbf{sparse}, distances between points become large and similar, and the amount of data needed to generalise grows \textbf{exponentially} with the number of dimensions.
\end{dfn}
\begin{itemize}
  \item In a unit square, a random point is near the border (within 0.001) with probability $0.4\%$; in 10\,000 dimensions $>99.999999\%$. In general $P=1-(1-2\delta)^d$.
  \item Average distance between two random points: $0.52$ in 2-D, $0.66$ in 3-D, about $408.25$ in 1\,000\,000-D $\Rightarrow$ new instances are far from training instances.
  \item To cover a space with $b$ bins per feature you need $b^d$ cells (e.g.\ $10^{100}$ for 100 features).
  \item Effect: high risk of \textbf{overfitting}; distance-based methods (k-NN, K-Means) suffer most.
  \item Remedy: more data (impractical: exponential) or \textbf{reduce dimensions}.
\end{itemize}
\begin{numer}[: near the border]
$\delta=0.001$. $d=2$: $1-0.998^2=0.004$ (0.4\%). $\ d=1000$: $1-0.998^{1000}=1-e^{-2}=\mathbf{0.865}$ (86\%).
\end{numer}

\section{Main Approaches}
\begin{dfn}{Projection}
In most real problems the training instances lie close to a much \textbf{lower-dimensional subspace}. \textbf{Projection} drops perpendicular onto that subspace (e.g.\ 3-D points near a plane are projected to 2-D). PCA is a projection method.
\end{dfn}
\begin{dfn}{Manifold learning}
A $d$-dimensional \textbf{manifold} is a shape that locally looks like a $d$-dimensional plane but is bent or twisted in a higher-dimensional space (e.g.\ the \textbf{Swiss roll}: 2-D manifold in 3-D). \textbf{Manifold learning} models this shape and \textbf{unrolls} it; projection would squash its layers together. Examples: LLE, Isomap, t-SNE.
\end{dfn}
\begin{itemize}
  \item \textbf{Manifold hypothesis:} most real high-dimensional data lie close to a much lower-dimensional manifold.
  \item \textbf{Feature selection} keeps a subset of original features; \textbf{feature extraction} creates new combined features (PCA).
  \item Linear methods: PCA, LDA. Non-linear: Kernel PCA, LLE, t-SNE, autoencoders.
\end{itemize}
\begin{drawit}{Projection vs.\ manifold (Swiss roll)}
\centering
\begin{tikzpicture}[scale=0.8]
\draw[thick,fill=mlblue!10] (0,0)--(3,0.5)--(3.8,2)--(0.8,1.5)--cycle; \node[font=\scriptsize] at (1.9,1) {plane};
\foreach \x/\y in {0.7/1.3,1.5/0.6,2.4/1.6,3/1.1,1.2/1.0,2.0/1.2}{\fill[mlorange] (\x,\y+0.35) circle(2pt); \draw[dotted] (\x,\y+0.35)--(\x,\y);}
\node[font=\scriptsize] at (1.9,-0.5) {projection onto a plane};
\begin{scope}[xshift=7cm,yshift=1cm]
\draw[very thick,mlpurple,domain=0:900,samples=200,smooth] plot ({0.0022*\x*cos(\x)},{0.0022*\x*sin(\x)});
\node[font=\scriptsize] at (0,-2.5) {Swiss roll: unroll, do not squash};
\end{scope}
\end{tikzpicture}
\end{drawit}

\section{Principal Component Analysis (PCA)}
\begin{dfn}{PCA}
\textbf{Principal Component Analysis} is an unsupervised, linear dimensionality-reduction technique that finds new orthogonal axes, the \textbf{principal components}, ordered so that the first captures the \textbf{maximum variance} of the data, the second the maximum remaining variance, and so on. The data is then projected onto the first $m$ components.
\end{dfn}
\begin{itemize}
  \item \textbf{Preserving variance:} the axis that keeps maximum variance also minimises the mean squared distance between original points and their projections (least information lost).
  \item Components are the \textbf{eigenvectors} of the covariance matrix; their \textbf{eigenvalues} are the variances along them.
  \item PCA assumes data is \textbf{centred}; scale features first if units differ.
\end{itemize}
\textbf{Algorithm:}
\begin{enumerate}
  \item Standardise / centre the data: $\tilde{\mathbf X}=\mathbf X-\bar{\mathbf x}$.
  \item Compute the covariance matrix $\mathbf C=\frac1{n-1}\tilde{\mathbf X}^{\!\top}\tilde{\mathbf X}$.
  \item Find eigenvalues and eigenvectors: $\mathbf C\vect w=\lambda\vect w$; sort by $\lambda$ (descending).
  \item Choose $m$ components (e.g.\ 95\% explained variance).
  \item Project: $\mathbf Z=\tilde{\mathbf X}\mathbf W_m$.
\end{enumerate}
In practice PCA is computed with \textbf{SVD}: $\tilde{\mathbf X}=\mathbf U\mathbf S\mathbf V^{\!\top}$; rows of $\mathbf V^{\!\top}$ are the components, $\lambda_j=s_j^2/(n-1)$.

\begin{numer}[: PCA by hand]
Points $(1,1),(2,3),(3,2),(4,5),(5,4)$. Means $\bar x=\bar y=3$; centred: $(-2,-2),(-1,0),(0,-1),(1,2),(2,1)$.\\
$s_{xx}=\frac{10}4=2.5$, $s_{yy}=2.5$, $s_{xy}=\frac{4+0+0+2+2}4=2$ $\Rightarrow$ $\mathbf C=\left[\begin{smallmatrix}2.5&2\\2&2.5\end{smallmatrix}\right]$.\\
$\det(\mathbf C-\lambda\mathbf I)=(2.5-\lambda)^2-4=0\Rightarrow\lambda_1=4.5$, $\lambda_2=0.5$.\\
$\vect w_1=(0.707,0.707)$, $\vect w_2=(0.707,-0.707)$. Explained variance: PC1 $\frac{4.5}{5}=\mathbf{90\%}$, PC2 $10\%$.\\
Project A: $z_1=\frac{-2-2}{\sqrt2}=-2.83$. Keeping only PC1 reduces 2-D to 1-D while keeping 90\% of the variance.
\end{numer}

\begin{drawit}{Principal components of 2-D data}
\centering
\begin{tikzpicture}[scale=0.7]
\draw[-Stealth] (0,0)--(6,0) node[below]{$x_1$}; \draw[-Stealth] (0,0)--(0,6) node[left]{$x_2$};
\foreach \x/\y in {1/1,2/3,3/2,4/5,5/4,1.5/2,3.5/3.8,2.5/2.2,4.2/4.4}{\fill[mlblue] (\x,\y) circle(3pt);}
\draw[very thick,mlred,-Stealth] (3,3)--(5.2,5.2) node[right,font=\scriptsize]{PC1 (90\%)};
\draw[very thick,mlgreen,-Stealth] (3,3)--(3.9,2.1) node[right,font=\scriptsize]{PC2 (10\%)};
\end{tikzpicture}
\end{drawit}

\begin{itemize}
  \item \textbf{Explained variance ratio} $=\lambda_j/\sum\lambda_i$. Choose $m$ so cumulative ratio $\ge95\%$, or look for the elbow of the \textbf{scree plot}. (MNIST: 784 $\to$ 154 dims keeps 95\%.)
  \item \textbf{Reconstruction:} $\hat{\mathbf X}=\mathbf Z\mathbf W_m^{\!\top}+\bar{\mathbf x}$; reconstruction error $=$ sum of discarded eigenvalues (compression).
  \item \textbf{Randomized PCA} (fast approximation), \textbf{Incremental PCA} (mini-batches, large data), \textbf{Kernel PCA} (non-linear via kernel trick).
\end{itemize}
\apa{reduces dimensions and noise, speeds up learning, removes correlation, helps visualisation, unsupervised}{components hard to interpret, linear only, loses some information, sensitive to scaling and outliers}{image compression, face recognition (eigenfaces), visualisation, gene data, preprocessing for ML}

\begin{center}
\small
\begin{tabular}{@{}p{3cm}p{6.6cm}p{6.6cm}@{}}
\toprule
 & \textbf{Feature selection} & \textbf{Feature extraction (PCA)}\\
\midrule
Output & subset of original features & new features (combinations)\\
Interpretability & high & low\\
Example & drop irrelevant columns & PCA, LDA\\
\bottomrule
\end{tabular}
\end{center}

\begin{imp}
Most asked: curse of dimensionality; projection vs.\ manifold learning; PCA steps with numerical (covariance, eigenvalues, explained variance); choosing number of components.
\end{imp}

\begin{iq}
\item \twom Define dimensionality reduction. Give two reasons for using it.
\item \twom What is the curse of dimensionality?
\item \twom What is a manifold?
\item \twom What is the explained variance ratio?
\item \fivem Explain projection and manifold learning.
\item \fivem How do you choose the number of principal components?
\item \fivem Compare feature selection and feature extraction.
\item \tenm Explain PCA with its algorithm and a numerical example.
\item \tenm Explain the curse of dimensionality and the main approaches to dimensionality reduction.
\end{iq}
